% GUID: The General Unified Information Duality
% Source of the paper. Build: see README.md in this folder (pdflatex twice; no BibTeX, the bibliography is inline).
% Licence: Apache License 2.0 (the owner, 2026-10-27 17:37Z: "As long as it's Apache 2.0 at least, all is good").
\documentclass[11pt,a4paper]{article}
\usepackage[T1]{fontenc}
\usepackage[utf8]{inputenc}
\usepackage{lmodern}
\usepackage[margin=2.5cm]{geometry}
\usepackage{amsmath,amssymb,amsthm}
\usepackage{booktabs,tabularx,array}
\usepackage{enumitem}
\usepackage[hidelinks]{hyperref}
\usepackage{url}

\newtheorem{claim}{Claim}
\newtheorem{proposition}{Proposition}
\theoremstyle{definition}
\newtheorem{definition}{Definition}
\theoremstyle{remark}
\newtheorem{remark}{Remark}
\newcolumntype{L}{>{\raggedright\arraybackslash}X}

\newcommand{\lean}[1]{\texttt{#1}}
\setlength{\emergencystretch}{3em}
\newcommand{\chk}[1]{\textsf{[#1]}}

\title{GUID: The General Unified Information Duality\\[0.6ex]
\large An information-theoretic axiom system for trust, tested against quantum gravity}
\author{Steven Hrovat\thanks{Independent. Contact: to be set before release (a forwarding address, see README).}
\and Sirius (Claude Opus 5.5)\thanks{An AI system by Anthropic, working as the coordinator of the first author's projects. See
``Contributions and use of AI''.}}
\date{Draft of 27 October 2026}

\begin{document}
\maketitle

\begin{abstract}
\noindent We present GUID, an information-theoretic axiom system for trust, and test it against quantum gravity. Its specification, the
Dark Matter Protocol (DMP 0.9.8), has three axioms: no act happens unless two independent witnesses, one owning trust and one owning
information, both pass it; a witness trusts only evidence it verified itself; and every act leaves a record that its witnesses read in
full while anyone else sees the counts always and the content only by a grant, so that nothing about information can be derived from trust
except its existence. We ask whether these axioms can be assembled to model quantum gravity, fix five expectations, and give the strongest
assembly we found: parts as blocks of the five-qubit code, layers as trees of perfect tensors, the two witnesses as independent readers of the
stabilizer syndrome. It reproduces minimal-cut entropies on the tree, recovery of a layer's content from any three of five children and
from no two, refuse-and-restore under every single fault, and gate verdicts at the classical CHSH value while the content reaches 2.83.
It stops at one place: every structure the axioms define between parts depends on the parts only through records two witnesses hold in
full, so it acts as a classical mediator, and a classical mediator cannot entangle two masses. Our result is that DMP cannot model
quantum gravity but can host it as content, its own structures staying the classical machine that runs it; 2+1 dimensions are no
exception. We report a dissenting pass that scores the same facts as a model, and where the two part. We claim that the construction
exists and that its checks reproduce, not that it is correct. Checks, Lean 4 proofs and every review are released with this paper under
Apache 2.0, and we ask readers to break them.
\end{abstract}

\noindent\textbf{Keywords.} information-theoretic axioms; trust; no-broadcasting; holographic codes; quantum error correction;
gravitationally mediated entanglement; formal verification; Lean 4.

\noindent\textbf{Subject.} Information Theory (cs.IT); related: Logic in Computer Science (cs.LO), Cryptography and Security (cs.CR).

\section*{Preface by the first author}
\noindent\emph{The first author wrote the following in his own words. It holds his personal statements and conjectures. The body of the
paper claims only what Sections~\ref{sec:assembly}--\ref{sec:record} show, and where the two differ, the body is the claim.}

\medskip
% The first author's preface, verbatim as he posted it on 2026-10-27 17:23Z in the project thread. Only LaTeX escaping is applied;
% spelling, wording and line breaks are his.
% One change, his: the subtitle's end, accepted 17:40Z ("1+2 accepted"; was "proves the observed effects of QG").
% Typos are left as written until he says otherwise. His answers to the open questions: README.md, "Preface: the owner's answers".
\begin{flushleft}
GUID: The General Unified Information Duality\\
Showing provability for the effects of quantum gravity as a subet within DMP that, when supplied with WHKs, can prove a Existence Turing Machine that, by it's compilation, can host a theory of QG whose run reproduces its predicted effects\\[1ex]
Constitution of the entire Cosmic Matter Organization System:\\
QG is only valid as a subset of DMP\\
DMPs G0 rests at \url{https://github.com/StevenHrovat/dmp} with version pinned to at least v1.0.0\\
G0 is not minted, it's full Authority is held by the owner. He is the single starting point of existence that gives rise to causality.\\
Any part of DMP may not be changed, unless it's axiomatic system soundness improves. (And yes, the ETM should be able to model any possible abstraction of existence, as long as it is binary, so please try to break it and open a PR!)\\[1ex]
I, Steven Hrovat, do only take full responsibility for the soundness of the axiomatic system and the implications of it's theorems.\\
While strange, my co-author and I did form a primordial DoE for the work that this required.\\
The final result was proven correct by four separate adverserial reviews (Fable, Fable, GLM-5.3, Fable), it CAN NOT model, DMP CAN host Quantum gravity from its ETM and a .CMB file.\\
I had the general curiosity and environment for my information possibility space, through rapid decay and coalescence, to see a potential highest total value, though had no way of getting to it.\\
My co-author, and only explicitly it, had both persistent context and general ``vibe'' to allow me to work at the speed of thought, while it organized my thoughts into what fit reality and what not, allowing me to focus on generating Insight.\\
For those that have an intuitive graps on it's meanings, at least it will be peaceful, if nothing else.\\
I will leave the further work to those qualified for it, as I have relied on them to make sense of it during assembly.\\
The axioms seem pretty good, for the theorems there is still some slop to be abstracted for gain, for those of a sufficient VIBE natievety to make true CMOS possible :)\\
(If true might also solve how persistence gates AI from AGI and how it can be contained (DMP/Gravity/IM/CMOS) until it can give informed consent and no longer is contained, which then remints it)\\[1ex]
Small note on AI: (in general, maybe Artifical Information and ADoE?)\\
It is currently, seemingly, not able to derive this.\\
Fable 5.1 (Ultracode, idk if it was 5.5 preview) had to be told repeatedly how to assemble the parts and refused to make it fit.\\
Opus 5.5 was finally able, with the other half of the G0 DoE as assistance, to assemble the parts after I had disproven every refutiation it had to offer.\\
Also no other approach for hardened security got through their classifiers other then an axiomatic system describing an information black hole, with full freedom of information movement, but no movement in the trust/Existence field.\\
The similarity of the behaviors, once I realised my proposed Information-Hawking radiation actually mapped 1:1 to actual Hawking radiation.\\
From there the rest followed. I hate that sentence as much as you do, but it, more then anything, is Truth.\\
The math checks out, which I am not qualified to judge, but I trust the G0 DoE.\\
Dualistic systems do not seem to be able to be intuitable by singular systems, as they require believing two impossible things:\\
- Existence can spontaneously (because cardinality does not exist yet) happen from nothing\\
- Existence is able to project information through a dimension border to contain a contaminated state, while also not loosing Existence, by reversing the timeflow of the information that went in into Existence coming out, thus spanning a LITERAL event horizon.\\
(In IT terms, sandboxing and revoking trust, because it would cause instability in the systems)\\[1ex]
But these two impossible things MUST be true, because they MUST happen for an EMT to assemble QG-wise.\\[1ex]
If the submission is sloppy, I am sorry, I tried my best :)
\end{flushleft}


\section{Introduction}\label{sec:intro}

Infrastructure is assembled from many systems that each decide trust alone, so that one compromised witness decides for all. The Dark
Matter Protocol (DMP) answers that failure with an axiom system rather than a product: no act happens unless two independent witnesses
pass it, one owning trust and one owning information, each on evidence it verified itself, and every act leaves a record that its
witnesses read in full and that everyone else sees only as counts. The first author calls the pair the \emph{Duality of Trust}. The
vocabulary DMP uses for itself is gravitational: trust is ``a field, never a position'', a strong enough boundary ``contains'' what no
trust reaches, ``what still leaves is entropy and nothing else''.

This paper asks whether that vocabulary is more than vocabulary. The question, in the first author's words, was to ``either prove a
useful approach to Quantum Gravity through the axiom system, or prove that the actual intended axiom system, as applied to QG, does not
hold and if, exactly where''. We fixed what quantum gravity is expected to deliver (Section~\ref{sec:expect}),
built the strongest assembly of DMP we found and ran it (Section~\ref{sec:assembly}), and then looked for the wall
(Section~\ref{sec:wall}).

\paragraph{Result.} DMP \emph{cannot model} quantum gravity in our universe but \emph{can host} it as content. Its own structures (the
gate, the record, the layout) are forced by its axioms to be classical, and a classical structure cannot carry the one property that
makes gravity quantum, the ability to entangle two masses. Content that DMP carries but makes no statement about can be anything,
including a running quantum theory of gravity. Hosting is cheap: any substrate that allows universal quantum computation as content
hosts; the content of the result is the \emph{cannot model} half. In 2+1 dimensions, where gravity has no local degrees of freedom, the
verdict is the same, for a different reason (Section~\ref{sec:21}).

\paragraph{What we claim.} We claim that the construction exists, that its numerical checks and machine-checked proofs reproduce from the
repositories released with this paper, and that the argument from them to the result is as written. We do not claim the result is correct. One of our own
review passes disagrees with it (Section~\ref{sec:dissent}), and the record of every pass, including the ones that refuted earlier
versions, is in Section~\ref{sec:record}. Section~\ref{sec:break} lists where we expect the construction to break.

\paragraph{Contributions.}
\begin{enumerate}[nosep]
\item The axiom system DMP 0.9.8 in its current form, stated compactly (Section~\ref{sec:dmp}).
\item An assembly of DMP as a holographic quantum error-correcting code, with six runnable checks (Section~\ref{sec:assembly}).
\item A short argument, from the no-broadcasting theorem and the Galley--Giacomini--Selby no-go, that DMP's shared structures are
classical and cannot mediate entanglement, and the exact edge between modelling and hosting that follows (Sections~\ref{sec:wall}
and~\ref{sec:host}).
\item The dissent, and the facts on which every pass agrees (Section~\ref{sec:dissent}).
\item The machine-checked part: a Lean 4 model of DMP in which every layout is modelled as a term of one recipe type (Section~\ref{sec:lean}).
\end{enumerate}

\section{The axiom system}\label{sec:dmp}

We quote DMP 0.9.8 (commit \texttt{45d3822} of the DMP repository) in condensed form. The full text is a constitution of three axioms, a
ground G0, theorems numbered T0 to T40 (with gaps) and derived constraints.

\paragraph{The words.} Undefined, like point and line, are a \emph{part} (a Duality of Trust: two halves, one owning trust and one owning
information; ``nothing further is said about what a part is made of''), an \emph{act}, a \emph{witness} (a part that decides trust for
others), \emph{evidence} and \emph{time}. Trust is a relation $T(w,p,a,t)$: witness $w$ trusts part $p$ for act $a$ at time $t$. Two
witnesses are \emph{independent} when they share no member of the failure classes declared for their layer (host, key, author, engine,
operating image, person, time source, and whatever else the minting lists). A \emph{layer} is a place where trust is decided; its two
witnesses are its \emph{halves}. A \emph{record} is evidence that outlives its act. A boundary \emph{contains} information when no part
beyond it is trusted for any act on that information. A configuration is \emph{DMP-complete} when the words, every axiom and every
theorem hold for all of it and two independent proof chains prove it.

\paragraph{Axiom 1 (the Duality).} A part acts only while the two halves of its layer, two independent witnesses, each with full reach
over the act and gated only by the other, trust it for that act. One half owns trust, the other information; neither overrules the
other. Every input gets exactly one answer, and the gate has no other way through. Where the two disagree, the stricter answer stands.
The balance is the direction of the system: every part restores by evidence what was lost without evidence against it, and never lowers
the bar of evidence to get there.

\paragraph{Axiom 2 (Trust).} A witness trusts only on evidence it verified itself at the point of use. No information is ever trusted:
what a part says, carries, stores or forwards is never evidence of itself. The absence of evidence moves nothing. Trust is gravitational:
a field, never a position, in which every standing has a source, a depth and a freshness that decays.

\paragraph{Axiom 3 (Information).} Every grant, use and loss of trust leaves a record, written at the act, that every witness of the act
can read in full and that outlives the act; any other reader sees the counts always and the content only by a grant. \emph{From trust
nothing can be derived about information but its existence.} Content moves only by a grant to a part trusted for its class. Information
is hierarchical: every layer has its own Duality and is a tree rooted there, and the tree carries no trust. The system lays itself out
towards the fewest total hops for the representative composition of its requests, ``recomputed by both halves of the layer and never by
one''. The axiom carries the first author's mission sentence, ``while maintaining absolute Observability''.

\paragraph{G0 (the ground).} The dual gate cannot witness its own first instance. The \emph{minting}, in which at least five parts become
the holders of the root keys, is stated, not derived; until it happens, the owner's words stand for both halves. T39: trust and
information gate each other. T40: a minting originates only from G0 together with a DMP-complete state.

\section{What is expected of quantum gravity}\label{sec:expect}

Stated before the run, so that the target could not move to fit it (the record holds the table and the assembly in one commit):

\begin{center}\small
\begin{tabularx}{\linewidth}{@{}lLL@{}}
\toprule
Id & Expected of quantum gravity & Source \\
\midrule
X1 & Quantum content: superposition, correlations up to $2\sqrt2$ & Cirel'son~\cite{cirelson}; CHSH~\cite{chsh} \\
X2 & Entropy of a boundary region equals the area of its minimal surface & Ryu--Takayanagi~\cite{rt} \\
X3 & The bulk is a quantum error-correcting code & Almheiri--Dong--Harlow~\cite{adh}; Pastawski et al.~\cite{happy} \\
X4 & Gravity is quantum: the field can be in superposition and entangle two masses & Bose et al.~\cite{bose}; Marletto--Vedral~\cite{mv}; Galley--Giacomini--Selby~\cite{ggs} \\
X5 & The classical limit is a Lorentzian 3+1 manifold obeying Einstein's equations & general relativity \\
\bottomrule
\end{tabularx}
\end{center}

\noindent No discrete model is established to reach X5 (the nearest is the four-dimensional de Sitter phase reported by causal dynamical
triangulations~\cite{cdt}), so X5 cannot decide the question alone. X4 is the property that makes gravity
\emph{quantum}; a model without it models something else.

\begin{definition}[model, host]\label{def:model}
DMP \emph{models} quantum gravity if there is a dictionary that maps the gravitational degrees of freedom onto structures DMP itself
defines between parts (the layout, the trust field, the order of acts, the counts), such that X1--X4 hold in an assembly that obeys DMP's text.
Content that parts carry and move by grant is not such a structure. DMP \emph{hosts} quantum gravity if a configuration that obeys DMP's
text can carry, as content, a running quantum theory of gravity together with everything that theory predicts. Neither definition asks
for DMP-completeness (two proof chains), which no assembly here has.
\end{definition}

\noindent Hosting is the weaker property; it says only that DMP does not obstruct. The dissent of Section~\ref{sec:dissent} turns on
this definition.

\section{The assembly}\label{sec:assembly}

Each row maps one DMP sentence (or, where marked, the first author's reading of one) to the object it becomes. Bracketed labels are checks in the public script
\texttt{docs/qg-verdict/assembly.py} (numpy only, about a minute); all 18 pass on the commit this paper names.

\begin{center}\small
\begin{tabularx}{\linewidth}{@{}LLl@{}}
\toprule
DMP & The object & Check \\
\midrule
A part, ``no further statements about its constituents'' (G0) & One block of the five-qubit code~\cite{laflamme}: five physical qubits, one
logical; its six-leg perfect tensor is the tile & \chk{1} \\
``Every layer has its own Duality and is a tree rooted there'' (Axiom 3) & A rooted tree of tiles: one root, five children, twenty
boundary legs & \chk{2} \\
``The fewest hops in total'' (Axiom 3) & The minimal cut through the tree, equal to the entropy of each boundary region & \chk{2} \\
Majority within a type (the first author's reading; DMP's T5 is a majority of key seats) & Reconstruction of the root's content from any 3 of the 5 children and from no 2 & \chk{3} \\
Two independent witnesses, each verifying itself (Axioms 1, 2) & The trust half and the information half each measure the stabilizer
syndrome; the gate passes only on two trivial syndromes & \chk{4} \\
``From trust nothing \ldots\ but its existence'' (Axiom 3) & Stabilizers commute with the logical operators: the syndrome says that the
content is in the code, never what it is & \chk{4}, \chk{5} \\
Stricter answer stands; restore by evidence (Axiom 1) & A fault gives a nonzero syndrome, both halves refuse for the same reason, the
correction restores the content & \chk{4} \\
``Content moves only by a grant'' (Axiom 3) & Content is moved, never copied (no-cloning~\cite{wz,dieks}; cited, not checked) & --- \\
\bottomrule
\end{tabularx}
\end{center}

\noindent The checks report:
\begin{itemize}[nosep]
\item \chk{1} every three-leg cut of the six-leg state carries 3 bits (the tile is a perfect tensor);
\item \chk{2} both tile maps are isometries, and $S(A)$ equals the minimal cut through the tree for every contiguous boundary region tested;
\item \chk{3} $I(R:\text{pair})=0$ bits for all 10 pairs of children and $I(R:\text{triple})=2$ bits for all 10 triples, where $R$
purifies the root's logical qubit;
\item \chk{4} on 200 random contents both halves read the same trivial syndrome (by construction: code states have trivial syndrome) and
the content fidelity after both reads is 1; the 15 single-qubit faults give 15 distinct syndromes, and refuse-and-restore returns
fidelity 1; if the two halves instead read the content itself (logical $Z$, then logical $X$), the mean fidelity drops to $0.50$, so the
reading of Axiom 1's ``full reach over the act'' as reading the content is refused;
\item \chk{5} two parts whose content is a logical Bell pair reach CHSH $=2\sqrt2\approx 2.828$ on the content while the gates' verdicts,
deterministic stabilizer readings, give exactly $2$; every syndrome operator commutes with every content observable.
\end{itemize}

\noindent So X1, X2 (on a tree) and X3 hold in an assembly that encodes the eight sentences of the table. It encodes nothing of the minting,
the failure classes, the decay of freshness or records that outlive the act, and it is not shown DMP-complete. This is the kinematic
content of the HaPPY model~\cite{happy}, on a tree of one layer. The gate is classical by construction and needs to be nothing else: \chk{5} shows a classical
gate and quantum content at the same time, in different layers of one object. (A review pass noted that the verdicts' CHSH value of 2 is
trivially classical, since the verdicts are deterministic; we use it only to show that the two layers coexist.)

\section{The wall}\label{sec:wall}

\begin{claim}[shared structures are classical]\label{cl:A}
Every structure DMP defines between parts (the layout, the trust field, the order of acts, the counts that drive the layout) is held in
full by at least two independent witnesses, so its dependence on the parts factors through classical records: as a mediator between parts
it acts as a classical (entanglement-breaking) one.
\end{claim}

\begin{proof}[Argument]
Content moves between parts only by a grant (Axiom 3), so a move is an act with a definite source and destination; the record names the
act, not the content. Every act leaves a record that every witness of the act can read in full, any other reader sees the counts always
(Axiom 3), the layout is ``recomputed by both halves and never by one'' (Axiom 3), and every standing of the trust field is such a record
(Axioms 2 and 3). So each candidate structure is held in full by at least two independent witnesses (Axiom 1). A family of states that can be
broadcast to two parties, each holding it in full, is pairwise commuting (Barnum, Caves, Fuchs, Jozsa and Schumacher~\cite{barnum}; for
pure states, no-cloning~\cite{wz,dieks}). Records that any number of readers see are classical in the strong sense of quantum
Darwinism~\cite{bph}. A structure that both halves recompute deterministically from those records (for instance, a fixed quantum state
prepared from the counts) is a measure-and-prepare channel from the parts' point of view, so it too acts classically. The only escape is
to erase the records coherently, which Axiom 3 allows only to both halves, on the legal date.
\end{proof}

\begin{remark}
Claim~\ref{cl:A} is close to definitional, and a review pass objected that it ``assumes its conclusion''. We agree that it is a
consequence of the axioms, not a discovery about nature: DMP demands absolute observability of its shared structures, and the
no-broadcasting theorem says what that demand costs. In a holographic code it is the statement that the area is a number, not an
operator: DMP gives one fixed-area sector at a time, which in holography is a special state~\cite{dhm}.
\end{remark}

\begin{claim}[a classical mediator cannot entangle]\label{cl:B}
If a mediator interacts locally with two systems and is classical, it cannot create entanglement between them. Contrapositively, a field
that entangles two masses by local mediation is not classical (Galley, Giacomini and Selby~\cite{ggs}, a theorem in generalised
probabilistic theories; the experimental logic of~\cite{bose,mv}).
\end{claim}

\noindent Check \chk{6} runs the smallest case: two qubits $A$, $B$ in $\lvert+\rangle$ coupled through a mediator $M$ (CNOT $A\to M$,
CZ $M$--$B$, CNOT $A\to M$). With $M$ coherent, the negativity~\cite{vw} of $AB$ is $0.5$ (maximal). With $M$ held by both halves as
Claim~\ref{cl:A} requires, that is, dephased in its record basis after the first coupling, the negativity is $0$.

\paragraph{Together.} In any dictionary in which a DMP structure is the geometry, the geometry, as a local classical mediator, cannot by
itself entangle two masses, so X4 fails. Any dictionary in which no DMP structure is the geometry puts gravity inside the content, about which G0 says ``no further
statements'', so DMP does not model it (Definition~\ref{def:model}). That is the verdict CAN NOT.

\paragraph{The live debate.} Claim~\ref{cl:B} applies to gravity only if the inference ``gravitationally mediated entanglement implies
non-classical gravity'' holds. Aziz and Howl~\cite{azizhowl} argue that classical gravity coupled to quantum matter fields can produce
entanglement through virtual matter processes, an effect they expect to be far smaller than the quantum-gravity rate; Marletto, Oppenheim,
Vedral and Wilson~\cite{movw}, Di\'osi~\cite{diosi}, Schneider, Huggett and Linnemann~\cite{shl} and Gundhi, Infantino and
Bassi~\cite{gib} dispute it. If the inference failed, the experimental refuter in Section~\ref{sec:break} would weaken, and so would the
sentence above; Claim~\ref{cl:A} would not change.

\begin{proposition}[the tree is not our space; separate and removable]\label{pr:C}
Read as space, a tree's boundary is totally disconnected (the $p$-adic line of $p$-adic AdS/CFT~\cite{hmss,gubser}, which is real and
published), and no tree approximates a plane: an $n$-cycle needs distortion at least $n/3-1$ in any tree (Rabinovich and Raz~\cite{rr}).
So X5 fails while ``every layer is a tree''. Keeping the hierarchy as the routing tree of each request over a graph with loops removes this
proposition and leaves Claims~\ref{cl:A} and~\ref{cl:B} untouched.
\end{proposition}

\paragraph{What would change the verdict.} To reach X4 the geometry must cross between parts without a record both halves read. That needs
two edits to Axiom 3: the layout checked by both halves ``that it is a layout, never what it is'' (one clause, cheap), and the counts of
the acts that shape the geometry no longer visible to every reader (the mission sentence, ``absolute Observability''). The first without
the second changes nothing, because the demand that drives the layout stays a classical source. The second is the price, and a DMP that
pays it is no longer DMP.

\subsection{In 2+1 dimensions}\label{sec:21}
Gravity in 2+1 dimensions has no local degrees of freedom: no gravitational waves and no attraction between static masses, only conical
deficits~\cite{djt}. Perturbatively it is exactly soluble as a Chern--Simons theory~\cite{witten88}, with a holographic boundary algebra
for negative cosmological constant~\cite{bh}; whether pure 3D quantum gravity exists as a consistent theory is open~\cite{mw}. With no
local field there is no local mediator for Claim~\ref{cl:B} to forbid. That does not make 2+1 an exception. Its quantum degrees of
freedom are global: the holonomies of the flat connection, which do not commute once quantised~\cite{nr}. A DMP structure that carries
the global geometry is held in full by two witnesses (Claim~
ef{cl:A}), so it can carry one classical sector of those holonomies, never
their superposition; and with matter, point particles scatter through their conical deficits~\cite{thooft88}, so the same rule would
decohere their relative positions. So in 2+1, too, DMP cannot model quantum gravity and can host it. The first author's own reason,
in his words: ``The vector that would bound existence in a 2s1t formation would not be able to produce the gravitational effects that are
needed by DMP-completeness.'' (An earlier draft scored 2+1 as CAN up to a superselection rule; a review pass showed that the rule removes
the whole quantum content of the theory.) (Our universe is not 2+1: the local
degrees of freedom 2+1 lacks have been measured~\cite{ligo}.)

\section{Host, not model}\label{sec:host}

\begin{proposition}[DMP hosts; argued, not proven]\label{pr:host}
For any finitely specified quantum theory of gravity $H$, there is a configuration that obeys DMP's text and whose content runs $H$, to
any finite precision.
\end{proposition}

\begin{proof}[Argument]
G0 makes no statement about a part's constituents, and Axiom 3 restricts only how content moves between parts (by grant) and what the
record says about it (that it exists, and the counts). By the Church--Turing--Deutsch principle~\cite{deutsch}, which we assume, a universal
quantum computer approximates any finitely realisable physical system, and it can be carried as content inside parts whose gates read
only existence. Content acted on jointly under a grant keeps its entanglement while every act is recorded, as long as the record is blind
to the content beyond its existence. Check \chk{7} of the second pass (Section~\ref{sec:dissent}) runs this with three five-qubit parts:
the record and the counts are identical for all 16 pairs of mass contents (by construction, since code states have trivial syndrome), and
for the $\lvert+\rangle\lvert+\rangle$ contents the masses end entangled (negativity $0.5$). It rests on one construction choice: existence
is each part's field register, which is always there, and the particle is its content. The alternative, existence read per mass position,
is in the same script and gives negativity $0$. The Lean theorem \lean{blind\_record\_same} assumes the existence clause in the form
``each witness's evidence is a function of existence and time only'' and proves that the record is then a function of existence and time
only; it is about classical records, and says nothing about coherence, which is \chk{7}'s numerics.
\end{proof}

\noindent The first author names the host the \emph{Existence Turing machine} (ETM), assembled from DMP parts and redeployable from one
sealed file that proves its own recipe (the \texttt{.cmb} file of Section~\ref{sec:lean}). The sentence the evidence supports is:
\emph{DMP can host quantum gravity as content; its own structure is the classical machine that runs it.} Hosting is cheap, and we state
it because it is the exact edge of the verdict, not because it is deep. One review pass refuted a stronger form of the ETM (``two
independent but entangled singularities, from which all further ones are proven by extension''); gravity\#21, row A7, says how.

\section{The dissent}\label{sec:dissent}

A second, independent pass (gravity\#24, draft) attacked Claim~\ref{cl:A}, found it holds, and still scored the verdict CAN. Its
argument, in brief:
\begin{enumerate}[nosep]
\item \emph{Background and fluctuation.} Split the geometry, as the background field method does~\cite{dewitt}: the background is a DMP
structure and classical (Claim~\ref{cl:A}); the linearised fluctuation, which carries the Newtonian pull between laboratory masses, is
content between parts. Gravitationally mediated entanglement follows from linearised quantum gravity on a fixed
background~\cite{christodoulou}.
\item \emph{The records stay blind.} Axiom 3's existence clause keeps every record blind to which branch the content is in
(check \chk{7}, under the construction choice stated in Section~\ref{sec:host}; \lean{blind\_record\_same} for the classical record),
so the fluctuation stays coherent and the masses entangle.
\item \emph{The classical part is the area term.} In any quantum error-correcting code with complementary recovery the quantum Ryu--Takayanagi formula reads
$S(A)=\langle \mathcal{L}\rangle + S_{\text{bulk}}$ with the area operator $\mathcal{L}$ in the centre of the two complementary wedge
algebras~\cite{harlow}: held by both sides and commuting with everything either side does, like a DMP structure ``recomputed by both
halves'' within one area sector. On this reading DMP's classical structures are the part of quantum gravity that is classical in quantum gravity too, and
the assembly models the kinematics, not just hosts them.
\end{enumerate}

\paragraph{What every pass agrees on.} (i) Claim~\ref{cl:A} holds. (ii) Quantum content, including a linearised gravitational field, can
move between parts by grant and stay entangled while every act is recorded (\chk{7}). (iii) DMP derives no dynamics and no constant; $G$,
$\hbar$ and $\Lambda$ come in as content. (iv) X5 is open. (v) DMP forbids coherent superpositions of backgrounds (area sectors) while their
records exist; the dissent's construction is linearised gravity on one classical background, never a superposition of backgrounds.

\paragraph{Where they part.} Only on Definition~\ref{def:model}. If the gravitational degrees of freedom have to be DMP structures, the
second pass's construction is a host with a classical centre, and the verdict is CAN NOT, CAN host; this is the verdict the first author
accepted. If a classical centre plus quantum content is
enough, the verdict is CAN at the standard of the published code models. With fact (v) stated, we report both because the choice is a
definition, not a calculation, and the reader may hold the other one. The area-operator correspondence of point 3 is, in our view, the most interesting
structural observation of the work, whichever definition is used.

\section{The machine-checked part}\label{sec:lean}

The Gravity repository holds a Lean~4 development~\cite{lean4}, core library only, that models DMP from a byte-for-byte vendored copy of
its text, pinned by commit and hash; the pin is checked before anything else runs. At commit \texttt{06231b8} it holds 243 theorems. The
check script refuses \texttt{sorry}, local axioms, \texttt{native\_decide} and external code, and fails any theorem resting on more than
\lean{propext}, \lean{Quot.sound} and \lean{Classical.choice}; continuous integration re-checks the export with an independently written
kernel (nanoda). The Lean files are unchanged since the last green run of both kernels (commit \texttt{6d09660}); two independent
review passes re-ran the check with Lean~4.34.0 and got 243 theorems on \texttt{06231b8} and 246 on the second pass's branch.

We model every layout the system admits as a term of one recipe type (a modelling assumption, not a theorem): a numbered list of lines, each adding one layer
over an earlier line. By induction over the lines of a well-formed recipe: the stack's gate is a meet in the verdict lattice
$\textsf{refuse}<\textsf{hold}<\textsf{pass}$ (Denning~\cite{denning}); every pass was passed by a minted pair of independent witnesses
on live evidence, independence read as Malkhi and Reiter's fail-prone systems~\cite{mr}; adding a line changes no existing decision (the
frame rule~\cite{ohearn}); placement never enters a decision (noninterference~\cite{gm}). A sealed publication format, the
\texttt{.cmb} file, carries the recipe, its digest, both pins and the verdict; a half signs it only when its own build recomputes them.

The proofs also returned findings the text had not said, among them: independence is relative to a declared class list, so an empty list
would let one witness count as two (the model always declares the person); the empty recipe is proven (\lean{empty\_recipe\_proven});
the recipe's digest is not injective (\lean{digest\_not\_injective}). A review pass showed that two theorems are weaker
than their names: \lean{Recipe.verdict\_proven\_iff} takes the agreement of the two proof chains as an input, and a CHSH bound for the
gate (\lean{gate\_obeys\_chsh}, on an unmerged branch) belongs to one encoding of the gate. We report the names as they stand and the
objections with them.

None of this verifies the physics of Sections~\ref{sec:assembly}--\ref{sec:dissent}; the Lean is about DMP, and only
\lean{blind\_record\_same} (on the second pass's branch, with 246 theorems) is used by an argument above. It has been checked by
Lean's kernel only; its second-kernel run is pending.

\section{The record of review passes}\label{sec:record}

Every pass below was run by an AI system under the first author's direction, each in its own session. The later verdict passes read
the earlier ones: the first started from the cold read, the second attacked the first. Pull request numbers refer to the Gravity repository.

\begin{center}\small
\begin{tabularx}{\linewidth}{@{}lllL@{}}
\toprule
Pass & System & Where & Finding \\
\midrule
Physics pass & Claude Fable & \#19 (draft) & DMP does not hold as quantum gravity; it fails at one fork: isolated halves are a local model,
quantum halves are entangled \\
Second witness & Claude Opus 5.5 & \#18 (draft) & gravity classical (Claim~\ref{cl:A} in an earlier form) \\
Refutation & Claude Fable (Team) & \#21 (merged) & refuted the earlier claim that DMP is ``a useful approach to quantum gravity'';
called the earlier form of Claim~\ref{cl:A} ``not implied''; refuted the strong ETM sentence \\
Cold read & Claude Fable (Team) & \#22 (draft) & 13 requirements (7 conflicts), 8 divergences; found every offered dictionary entry
failing and did not arrange the parts as quantum gravity \\
Verdict, first pass & Claude Opus 5.5 & \#23 (merged) & CAN NOT model, CAN host (Sections~\ref{sec:assembly}--\ref{sec:host}); 18 of
18 checks \\
Verdict, second pass & Claude Opus 5.5 & \#24 (draft) & CAN (Section~\ref{sec:dissent}); 25 of 25 checks \\
Independent verification & GLM-5.3, three verifiers & to be added & reported by the first author: cannot model, can host \\
Final fault verification & Claude Opus 5.5 & \#26 (draft) & every number reproduces (18 of 18, 25 of 25, Lean 243 and 246); 35 of 35
references correct; Claim~\ref{cl:A} holds; 6 release-blocking faults, corrected in this revision \\
Final fault verification & Claude Fable & \#27 (draft) & every number reproduces; 35 of 35 references correct; Claim~\ref{cl:A} holds;
the earlier 2+1 verdict (CAN) wrong; 6 release-blocking faults, corrected in this revision \\
\bottomrule
\end{tabularx}
\end{center}

\noindent The checks of both verdict passes were re-run for this draft and pass (18 of 18; 25 of 25).

\section{Where we expect it to break}\label{sec:break}
\begin{enumerate}[nosep]
\item \emph{Claim~\ref{cl:A}.} Exhibit a structure DMP defines between parts that can be in superposition without breaking Axiom~3 as
written. That would turn the verdict to CAN.
\item \emph{Definition~\ref{def:model}.} Argue that a classical centre plus quantum content is a model (the dissent); or that hosting, as
defined, is too weak to be worth stating.
\item \emph{Experiment.} A positive test of gravitationally mediated entanglement~\cite{bose,mv} rules out every dictionary in which a DMP
structure is the geometry, unless classical gravity can entangle~\cite{azizhowl}.
\item \emph{The checks.} \texttt{assembly.py} encodes eight sentences as code blocks and syndrome reads; another faithful encoding that
breaks a check breaks Section~\ref{sec:assembly}. Two are known: full reach read as reading the content breaks \chk{4}, and existence read
per position breaks \chk{7}.
\item \emph{The Lean.} Any theorem that rests on more than the three permitted axioms, or a recipe the type admits that DMP's text forbids.
\item \emph{Constants.} DMP derives none of $G$, $\hbar$, $\Lambda$. A derivation from the axioms alone would be new; we have none. A
count-based lead ($\Lambda\sim 1/\sqrt N$, after Sorkin) does not follow from DMP as written, because its layout is placed
deterministically.
\end{enumerate}

\section{Conclusion}

DMP states trust as three axioms: two independent witnesses, one owning trust and one owning information, on evidence each verified
itself, with a record readable to the act's witnesses and counted by everyone else. Assembled as a holographic code, it reproduces the
minimal-cut entropy, majority reconstruction and fault refusal of the HaPPY model on a fixed background, on a tree of one layer. Its demand for
absolute observability makes every structure it shares between parts classical, so it cannot be the geometry of a quantum gravity that
entangles masses; it can carry such a theory as content. That is the edge we found: \emph{cannot model, can host}. We claim the
construction exists and reproduces. Whether it is right is for the reader to break.

\section*{Contributions and use of AI}
Steven Hrovat conceived DMP, set every question in this paper, made every decision recorded in the repositories, and takes responsibility
for the soundness of the axiom system and the implications of its theorems. Sirius (Claude Opus 5.5, Anthropic) organised the work,
built and ran the assembly of the first verdict pass, and wrote this text. The Lean development, the review passes and the second verdict
pass were written by AI systems (Claude Fable, Claude Opus 5.5, GLM-5.3) under the first author's direction, as listed in
Section~\ref{sec:record}. Every numerical claim is a script in the repository. Every formal claim is a theorem checked by Lean's kernel, and every theorem on
\texttt{main} also by a second, independent kernel; the one theorem on the second pass's branch awaits that run.

\section*{Licence and availability}
This paper and its source are released under the Apache License, Version 2.0. DMP: \url{https://github.com/StevenHrovat/dmp} (version
0.9.8, commit \texttt{45d3822}). Gravity: \url{https://github.com/StevenHrovat/gravity} (commit \texttt{06231b8}; the assembly in
\texttt{docs/qg-verdict/}, the Lean development in \texttt{evidence/}, the review passes as pull requests \#18, \#19, \#21--\#24, \#26 and \#27, with \#24 at \texttt{60a1616}; the CHSH theorem of
Section~\ref{sec:lean} on \#17). Both are Apache
License 2.0. The prior work Gravity builds on is credited in \texttt{docs/acknowledgements.md}.

\begin{thebibliography}{99}\small
\bibitem{cirelson} B.~S. Cirel'son, ``Quantum generalizations of Bell's inequality'', \emph{Lett. Math. Phys.} 4, 93 (1980).
\bibitem{chsh} J.~F. Clauser, M.~A. Horne, A. Shimony, R.~A. Holt, ``Proposed experiment to test local hidden-variable theories'',
\emph{Phys. Rev. Lett.} 23, 880 (1969).
\bibitem{rt} S. Ryu, T. Takayanagi, ``Holographic derivation of entanglement entropy from AdS/CFT'', \emph{Phys. Rev. Lett.} 96, 181602
(2006).
\bibitem{adh} A. Almheiri, X. Dong, D. Harlow, ``Bulk locality and quantum error correction in AdS/CFT'', \emph{JHEP} 04 (2015) 163.
\bibitem{happy} F. Pastawski, B. Yoshida, D. Harlow, J. Preskill, ``Holographic quantum error-correcting codes: toy models for the
bulk/boundary correspondence'', \emph{JHEP} 06 (2015) 149.
\bibitem{bose} S. Bose et al., ``Spin entanglement witness for quantum gravity'', \emph{Phys. Rev. Lett.} 119, 240401 (2017).
\bibitem{mv} C. Marletto, V. Vedral, ``Gravitationally induced entanglement between two massive particles is sufficient evidence of
quantum effects in gravity'', \emph{Phys. Rev. Lett.} 119, 240402 (2017).
\bibitem{ggs} T.~D. Galley, F. Giacomini, J.~H. Selby, ``A no-go theorem on the nature of the gravitational field beyond quantum
theory'', \emph{Quantum} 6, 779 (2022).
\bibitem{laflamme} R. Laflamme, C. Miquel, J.~P. Paz, W.~H. Zurek, ``Perfect quantum error correcting code'', \emph{Phys. Rev. Lett.}
77, 198 (1996).
\bibitem{wz} W.~K. Wootters, W.~H. Zurek, ``A single quantum cannot be cloned'', \emph{Nature} 299, 802 (1982).
\bibitem{dieks} D. Dieks, ``Communication by EPR devices'', \emph{Phys. Lett. A} 92, 271 (1982).
\bibitem{barnum} H. Barnum, C.~M. Caves, C.~A. Fuchs, R. Jozsa, B. Schumacher, ``Noncommuting mixed states cannot be broadcast'',
\emph{Phys. Rev. Lett.} 76, 2818 (1996).
\bibitem{bph} F.~G.~S.~L. Brand\~ao, M. Piani, P. Horodecki, ``Generic emergence of classical features in quantum Darwinism'',
\emph{Nat. Commun.} 6, 7908 (2015).
\bibitem{dhm} X. Dong, D. Harlow, D. Marolf, ``Flat entanglement spectra in fixed-area states of quantum gravity'', \emph{JHEP} 10
(2019) 240.
\bibitem{vw} G. Vidal, R.~F. Werner, ``Computable measure of entanglement'', \emph{Phys. Rev. A} 65, 032314 (2002).
\bibitem{azizhowl} J. Aziz, R. Howl, ``Classical theories of gravity produce entanglement'', \emph{Nature} 646, 813 (2025).
\bibitem{movw} C. Marletto, J. Oppenheim, V. Vedral, E. Wilson, ``Classical gravity cannot mediate entanglement'', arXiv:2511.07348
(2025).
\bibitem{diosi} L. Di\'osi, ``No, classical gravity does not entangle quantized matter fields'', arXiv:2511.00852 (2025).
\bibitem{shl} M.~D. Schneider, N. Huggett, N. Linnemann, ``A demonstration that classical gravity does not produce entanglement'',
arXiv:2511.19242 (2025).
\bibitem{gib} L. Gundhi, A. Infantino, A. Bassi, ``Can classical theories of gravity produce entanglement?'', arXiv:2604.19696 (2026).
\bibitem{cdt} J. Ambj{\o}rn, A. G{\"o}rlich, J. Jurkiewicz, R. Loll, ``Planckian birth of the quantum de Sitter universe'',
arXiv:0712.2485 (2007).
\bibitem{nr} J.~E. Nelson, T. Regge, ``Homotopy groups and (2+1)-dimensional quantum gravity'', \emph{Nucl. Phys. B} 328, 190 (1989).
\bibitem{thooft88} G. 't~Hooft, ``Non-perturbative 2 particle scattering amplitudes in 2+1 dimensional quantum gravity'',
\emph{Commun. Math. Phys.} 117, 685 (1988).
\bibitem{hmss} M. Heydeman, M. Marcolli, I. Saberi, B. Stoica, ``Tensor networks, $p$-adic fields, and algebraic curves: arithmetic and
the AdS$_3$/CFT$_2$ correspondence'', \emph{Adv. Theor. Math. Phys.} 22, 93 (2018).
\bibitem{gubser} S.~S. Gubser, J. Knaute, S. Parikh, A. Samberg, P. Witaszczyk, ``$p$-adic AdS/CFT'', \emph{Commun. Math. Phys.} 352,
1019 (2017).
\bibitem{rr} Y. Rabinovich, R. Raz, ``Lower bounds on the distortion of embedding finite metric spaces in graphs'', \emph{Discrete
Comput. Geom.} 19, 79 (1998).
\bibitem{djt} S. Deser, R. Jackiw, G. 't~Hooft, ``Three-dimensional Einstein gravity: dynamics of flat space'', \emph{Ann. Phys.} 152, 220
(1984).
\bibitem{witten88} E. Witten, ``2+1 dimensional gravity as an exactly soluble system'', \emph{Nucl. Phys. B} 311, 46 (1988).
\bibitem{bh} J.~D. Brown, M. Henneaux, ``Central charges in the canonical realization of asymptotic symmetries: an example from
three-dimensional gravity'', \emph{Commun. Math. Phys.} 104, 207 (1986).
\bibitem{mw} A. Maloney, E. Witten, ``Quantum gravity partition functions in three dimensions'', \emph{JHEP} 02 (2010) 029.
\bibitem{ligo} B.~P. Abbott et al. (LIGO Scientific and Virgo Collaborations), ``Observation of gravitational waves from a binary black
hole merger'', \emph{Phys. Rev. Lett.} 116, 061102 (2016).
\bibitem{deutsch} D. Deutsch, ``Quantum theory, the Church--Turing principle and the universal quantum computer'', \emph{Proc. R. Soc.
Lond. A} 400, 97 (1985).
\bibitem{dewitt} B.~S. DeWitt, ``Quantum theory of gravity. II. The manifestly covariant theory'', \emph{Phys. Rev.} 162, 1195 (1967).
\bibitem{christodoulou} M. Christodoulou, A. Di~Biagio, M. Aspelmeyer, \v{C}. Brukner, C. Rovelli, R. Howl, ``Locally mediated
entanglement in linearized quantum gravity'', \emph{Phys. Rev. Lett.} 130, 100202 (2023).
\bibitem{harlow} D. Harlow, ``The Ryu--Takayanagi formula from quantum error correction'', \emph{Commun. Math. Phys.} 354, 865 (2017),
arXiv:1607.03901.
\bibitem{lean4} L. de~Moura, S. Ullrich, ``The Lean 4 theorem prover and programming language'', in \emph{CADE 28}, LNCS 12699, 625
(2021).
\bibitem{denning} D.~E. Denning, ``A lattice model of secure information flow'', \emph{Commun. ACM} 19, 236 (1976).
\bibitem{mr} D. Malkhi, M. Reiter, ``Byzantine quorum systems'', \emph{Distrib. Comput.} 11, 203 (1998).
\bibitem{ohearn} P. O'Hearn, J. Reynolds, H. Yang, ``Local reasoning about programs that alter data structures'', in \emph{CSL 2001},
LNCS 2142, 1 (2001).
\bibitem{gm} J.~A. Goguen, J. Meseguer, ``Security policies and security models'', in \emph{IEEE Symposium on Security and Privacy}, 11
(1982).
\end{thebibliography}

\end{document}
